\begin{table}[htbp]
\centering
\footnotesize
\caption{\textbf{4.4} Ljung--Box sobre los residuos. En S2 se descartan los primeros 7 residuos, que son artefactos del arranque de las diferencias $(1-B)(1-B^6)$. Los grados de libertad bajan a $m-2$ en ambos modelos.}
\label{tab:p4_diagnostico}
\begin{tabular}{ccccccc}
\toprule
\textbf{modelo} & \textbf{$m$} & \textbf{$n$ residuos} & \textbf{$Q(m)$} & \textbf{gl} & \textbf{valor-$p$} & \textbf{ruido blanco} \\
\midrule
S1 & 6 & 245 & 9.72 & 4 & 0.0454 & no \\
S1 & 12 & 245 & 23.36 & 10 & 0.0095 & no \\
S2 & 6 & 231 & 4.31 & 4 & 0.3655 & \textbf{s\'i} \\
S2 & 12 & 231 & 11.76 & 10 & 0.3015 & \textbf{s\'i} \\
\bottomrule
\end{tabular}
\\[0.3em]{\scriptsize \textbf{Qu\'e leer:} S2 es el \'unico que deja residuos de ruido blanco. S1 es estacionario y supone una media fija que no existe, as\'i que la tendencia de agosto en adelante se le queda en los residuos. Este resultado anticipa correctamente qui\'en gana en la Parte~5.}
\end{table}
